\clearpage
\setcounter{chapter}{4}
\chaptermark{干涉装置与光场时空相干性\quad 激光}
\pdfbookmark[0]{第四章 干涉装置与光场时空相干性 激光}{chapter4}
\problem{4-1}{
一菲涅耳双面镜之夹角为 $20'$，一缝光源平行交棱且与交棱距离为 $10\,\mathrm{cm}$，接收屏幕在两个相干像光源的正前方，且与交棱相距 $210\,\mathrm{cm}$，设光波长为 $600\,\mathrm{nm}$。
\begin{enumerate}
\item 干涉条纹的间距 $\Delta x$ 为多少？
\item 在屏幕上最多能看到几个条纹？
\item 如果光源与交棱之距离维持不变，而在横向作一小位移 $\delta s$，幕上干涉条纹有何变化？
\item 如果计及缝光源的宽度，要求干涉场的衬比度不至于为零，则允许缝光源的最大宽度 $b_{\mathrm M}$ 为多少？
\end{enumerate}
}
\problem{4-2}{
一束平行光正入射于一双棱镜，其顶角为 $3.5'$，折射率为 $1.5$，相距 $5.0\,\mathrm m$ 处置放一屏幕。光波长为 $500\,\mathrm{nm}$。
\begin{enumerate}
\item 求幕上干涉条纹的间距 $\Delta x$。
\item 求幕上出现的条纹数目 $N$。
\end{enumerate}
}
\problem{4-3}{
一菲涅耳双棱镜，其顶角为 $1.5^\circ$，折射率为 $1.5$，相距点光源 $6.0\,\mathrm{cm}$，相距屏幕 $310\,\mathrm{cm}$，光波长为 $500\,\mathrm{nm}$。求幕上干涉条纹的间距 $\Delta x$。
}
\problem{4-4}{
一劳埃德镜其镜面宽度为 $5.0\,\mathrm{cm}$，一缝光源在镜面左侧，离边缘 $2.0\,\mathrm{cm}$，比镜面高出 $a=0.5\,\mathrm{mm}$；接收屏幕在镜面右侧，离边缘 $300\,\mathrm{cm}$。设光波长为 $589\,\mathrm{nm}$。
\begin{enumerate}
\item 求幕上条纹的间距 $\Delta x$，幕上出现条纹的数目 $N$。
\item 若缝光源平移从而改变了它离镜面的高度，则条纹将发生怎样的变化？这种变化与习题 4-1（3）双面镜情形有何不同？
\item 若计及缝光源宽度 $b$ 的影响，试定性画出幕上干涉强度 $I(x)$ 的变化曲线。\par
提示：回顾第 2 章 2-14 题，两者有可类比性。
\item 设光源宽度为 $b$，其下边距离镜面高度为 $a$，屏幕与两个相干光源的距离为 $D$，试证明，干涉强度函数由下式给出，
\[
I(x)=I_0\left(1+\gamma(x)\cos\left(2\pi\frac{2a+b}{\lambda D}x\right)\right),
\]
\[
\gamma(x)=\gamma_0\frac{\sin\left(2\pi\frac{b}{\lambda D}x\right)}{2\pi\frac{b}{\lambda D}x}.
\]
这里，$\gamma_0$ 是该劳埃德镜对单一点源所造成的干涉场的衬比度，它取决于镜面的光强反射率；若考虑到镜面的半波损，那只需在余弦函数中添加 $\pi/2$ 相移量，这不影响 $I(x)$ 曲线的变化特点。
\end{enumerate}
}
\problem{4-5}{
如图所示，这是一种利用干涉条纹的移动来测量气体折射率的原理性结构。在双缝之一 $S_1$ 后面置放一长度为 $l$ 的透明容器，当待测气体徐注入容器而将空气逐渐排出，在此过程中，观察者视场中的条纹就将移动。人们可由条纹移动的方向和数目，测定气体的折射率。
\begin{enumerate}
\item 若待测气体的折射率大于空气之折射率，试预测干涉条纹怎样移动？
\item 设 $l$ 为 $2.0\,\mathrm{cm}$，光波长为 $589.3\,\mathrm{nm}$，空气折射率 $n_0$ 为 $1.000276$，观测到条纹向上移动了 $20$ 个，求待测氢气的折射率 $n$。
\end{enumerate}
\tikzfig{fig_4_5}{4-5}
}
\problem{4-6}{
类似上题的干涉装置如图所示，称作瑞利干涉仪，用以测定空气折射率。在双缝后面置放两个透明长管 $T_1$ 和 $T_2$，其中 $T_2$ 管已充满待测的空气，而 $T_1$ 管在初始时刻为真空，然后徐徐注入空气，直至充满气压相同于 $T_2$ 管的空气。测定这一过程中观察点 $P$ 强度变化的次数 $N$ 为 $98.5$，$T_1$ 管中空气柱长度为 $20\,\mathrm{cm}$，光波长为 $589.3\,\mathrm{nm}$。试求出空气的折射率 $n_0$。
\tikzfig{fig_4_6}{4-6}
}
\problem{4-7}{
现用钠光灯作为杨氏双缝干涉实验的光源，其宽度 $b$ 已被一光阑限制为 $2\,\mathrm{mm}$，它与双缝平面相距 $2.5\,\mathrm m$。为了在幕上能出现可见的干涉条纹，问：
\begin{enumerate}
\item 双缝间隔不能大于多少？设其为 $d_0$。
\item 若双缝实际间距 $d$ 取 $d_0/2$，此时前方幕上干涉场的衬比度 $\gamma$ 为多少？
\end{enumerate}
}
\problem{4-8}{
一个直径为 $1\,\mathrm{cm}$ 的热光源，如果用干涉孔径角 $\Delta\theta_0$ 来描述，其空间相干范围为多少弧度？如果用相干面积 $\Delta S_0$ 来描述，问 $1\,\mathrm m$ 远的相干面积 $\Delta S_{10}$ 为多少？$10\,\mathrm m$ 远的相干面积 $\Delta S_{20}$ 为多少？取光波长为 $550\,\mathrm{nm}$。
}
\problem{4-9}{
若直接以月亮作为光源，在地面上作杨氏双孔实验，为了获得可见的干涉条纹，问双孔间隔 $d$ 不能大于多少？已知月地距离约为 $3.8\times10^5\,\mathrm{km}$，月球直径为 $3477\,\mathrm{km}$，光波长取 $550\,\mathrm{nm}$。
}
\problem{4-10}{
一直径为 $d$ 的细丝垫在两块平板玻璃之一边，以形成楔形空气层，在钠黄光垂直照射下出现干涉条纹如图所示。试求细丝直径 $d$。
\tikzfig{fig_4_10}{4-10}
}
\problem{4-11}{
块规是机械加工技术中所用的一种长度标准，它是一块钢质的长方体，其两个端面经过研磨抛光，达到相互平行，如图所示，是两个相同规号的块规，其中 $G_1$ 的长度是标准的，$G_2$ 是待校准的。校准的方法如下：把 $G_1$ 和 $G_2$ 放在钢质平台上并使之严密接触；再用一块透明平板 $T$ 压在 $G_1$、$G_2$ 上面，以形成一楔形空气层，如果 $G_1$ 和 $G_2$ 的高度略有差别，在单色光照射下便出现等厚干涉条纹，用以精测待校准块规长度的偏差。
\begin{enumerate}
\item 设光波长为 $5893\,\text{\AA}$，$G_1$、$G_2$ 相距 $l$ 为 $5\,\mathrm{cm}$，出现于 $G_1$ 区和 $G_2$ 区的条纹间距均为 $0.50\,\mathrm{mm}$，试求高度差 $\Delta h$。
\item 怎样判断 $G_1$、$G_2$ 高度谁高谁低？
\item 若出现两区中条纹间距不相等，比如 $G_1$ 区的 $\Delta x_1$ 为 $0.5\,\mathrm{mm}$，而 $G_2$ 区的 $\Delta x_2=0.3\,\mathrm{mm}$，这反映出什么问题？
\end{enumerate}
\tikzfig{fig_4_11}{4-11}
}
\problem{4-12}{
如图所示它是一种干涉膨胀计，其中 $G$ 为两个标准的石英环，$C$ 为待测的柱形样品，其略有倾斜的上表面与石英盖板 $T_1$ 之间形成一楔形空气层，从而产生等厚条纹。当温度改变时，由于样品与石英有不同的膨胀系数，以致空气层厚度发生了变化而出现条纹移动。
\begin{enumerate}
\item 设样品和石英环的高度 $l\approx1\,\mathrm{cm}$，在温度升高 $\Delta t\approx100\,^\circ\mathrm C$ 过程中，视场中的干涉条纹移动了 $20$ 条，光波长为 $5893\,\text{\AA}$。求出该样品的线膨胀系数与石英的差值 $\Delta\beta$。
\item 怎样判断差值 $\Delta\beta$ 的正负号？设石英膨胀系数为 $3.5\times10^{-7}/^\circ\mathrm C$，条纹移动方向背向交棱，试给出样品的线膨胀系数 $\beta$ 值。
\end{enumerate}
\tikzfig{fig_4_12}{4-12}
}
\problem{4-13}{
一待测透镜与玻璃平晶生成的牛顿环，从其中心往外数第 $5$ 环和第 $10$ 环的直径分别为 $1.4\,\mathrm{mm}$ 和 $3.4\,\mathrm{mm}$。求出透镜的曲率半径，设光波长为 $0.63\,\mu\mathrm m$。
}
\problem{4-14}{
现观察到一肥皂膜的反射光呈现绿色，这时视线与膜法线之夹角约为 $35^\circ$。
\begin{enumerate}
\item 试估算此膜的最小厚度 $h_{\mathrm m}$ 约为多少？设肥皂水的折射率为 $1.33$，绿光波长为 $515\,\mathrm{nm}$。
\item 若观察者的注视点不变动，而改变了视线角度，从 $35^\circ$ 开始逐渐变小，试问该处先后呈现何种色调？可能先后出现绿、黄绿、黄、橙、红这种色调变化吗？
\end{enumerate}
}
\problem{4-15}{
以玻璃片为衬底，涂上一层透明薄膜，其折射率为 $1.30$，设玻璃折射率为 $1.5$。
\begin{enumerate}
\item 对于波长为 $550\,\mathrm{nm}$ 的光而言，这膜厚应当取多少才能使其反射光因干涉而相消？这时光强反射率 $R$ 为多少？
\item 此膜厚对于波长为 $400\,\mathrm{nm}$ 的紫光或 $700\,\mathrm{nm}$ 的红光，其反射双光束之间分别有多大的相位差 $\delta_1$ 或 $\delta_2$。
\end{enumerate}
}
\problem{4-16}{
GaAs 发光管被制成半球形，以增加位于球心的发光区的对外输出功率。为了进一步提高输出功率，常在表面镀上一层增透膜如图所示。GaAs 发射光波长为 $930\,\mathrm{nm}$，折射率为 $3.4$。
\begin{enumerate}
\item 无增透膜时，球面光强反射率 $R_0$ 为多少？
\item 为了实现完全消反射，增透膜的折射率 $n_{\mathrm t}$ 和厚度 $h_{\mathrm t}$ 应当为多少？
\item 如果选用折射率为 $1.38$ 的氟化镁 $\mathrm{MgF_2}$ 能否增透？此时其光强反射率 $R'$ 为多少？膜层厚度 $h'$ 应取多少？
\item 若选用折射率为 $2.58$ 的硫化锌 $\mathrm{ZnS}$ 能否增透？此时其光强反射率 $R''$ 为多少？膜层厚度 $h''$ 应为多少？
\end{enumerate}
\tikzfig{fig_4_16}{4-16}
}
\problem{4-17}{
用眼睛或透镜观察等倾干涉环，里疏外密。试证明，在傍轴条件下，干涉环的半径 $\rho_m$ 与条纹级数差 $m$ 之关系为
\[
\rho_m\approx\sqrt m\cdot\rho_0,\qquad \rho_0=f\sqrt{\frac{n\lambda}{h}}.
\]
这里，整数 $m$ 为该干涉环与最靠近中心那干涉环之间的级数差，亦即两者之间的环数；$\rho_0$ 为一常数，$f$ 为焦距。
}
\problem{4-18}{
用钠光观察迈克耳孙干涉条纹，先看到干涉场中有 $12$ 个亮环，且中心是亮的；尔后移动一臂镜面 $M_1$，看到中心吞（吐）了 $10$ 环，而此时干涉场中还存在 $5$ 个亮环。试求：
\begin{enumerate}
\item 镜面 $M_1$ 移动的距离 $\Delta h$。
\item 开始时中心亮斑的干涉级 $k_0$，相应的等效空气膜厚 $h_0$。
\end{enumerate}
}
\problem{4-19}{
钠灯发射的黄光包含两条相近的谱线，俗称黄双线 $\lambda_1$ 和 $\lambda_2$。以钠黄光入射于迈克耳孙干涉仪，通过镜面移动来观察干涉场衬比度的变化，从而可分辨出黄双线的波长差 $\Delta\lambda$。
\begin{enumerate}
\item 实测结果为，在条纹由最清晰到最模糊过程中，视场里吞（吐）了 $490$ 个条纹，求钠双线的波长差 $\Delta\lambda$，以及 $\lambda_1$ 和 $\lambda_2$。已知事先粗测了钠黄光波长 $\bar\lambda$ 为 $5893\,\text{\AA}$。
\item 水银灯发射的光谱中也含有较强的黄双线，其波长为 $\lambda'_1\approx5770\,\text{\AA}$、$\lambda'_2\approx5791\,\text{\AA}$。在上述实验中若改用这水银黄光，则视场里由最清晰到最模糊过程中吞吐的条纹数 $N'$ 为多少？
\end{enumerate}
}
\problem{4-20}{
用迈克耳孙干涉仪进行精密测长，入射光为 $6328\,\text{\AA}$ 的 He-Ne 激光，其谱线宽度为 $10^{-3}\,\text{\AA}$，对于光强度讯号的测量灵敏度可达 $1/8$ 个条纹。
\begin{enumerate}
\item 这台干涉测长仪的测长精度 $\delta l$ 为多少？
\item 这台仪器一次测长量程 $l_{\mathrm M}$ 为多少？
\end{enumerate}
}
\problem{4-21}{
迈克耳孙干涉仪中的一臂镜面以速度 $v$ 匀速推移，而用透镜接收干涉条纹，并将它会聚到光电元件上，把光强变化转化为电讯号。
\begin{enumerate}
\item 若测得电讯号的时间频率为 $f\,(\mathrm{Hz})$，求入射光的波长 $\lambda$。
\item 若入射光波长在可见光谱区中间 $0.55\,\mu\mathrm m$ 左右，要使电讯号的频率控制在低频范围，比如 $50\,\mathrm{Hz}$，问反射镜平移速度应当为多少？
\item 若入射光波长在红外光谱区，比如 $20\,\mu\mathrm m$，要使电讯号的频率控制在低频范围，比如 $100\,\mathrm{Hz}$，问反射镜平移速度应当为多少？
\item 若反射镜平移速度为 $30\,\mu\mathrm{m/s}$，则钠黄光入射时所产生的电讯号其拍频 $f_{\mathrm b}$ 为多少？已知钠黄光双线 $\lambda_1=5890\,\text{\AA}$、$\lambda_2=5896\,\text{\AA}$。
\end{enumerate}
}
\problem{4-22}{
镉灯为一准单色光源，其发射的中心波长 $\lambda_0$ 为 $6428\,\text{\AA}$，谱线宽度 $\Delta\lambda$ 为 $10^{-2}\,\text{\AA}$。
\begin{enumerate}
\item 求其光场的相干长度 $L_0$ 和相干时间 $\tau_0$。
\item 求镉红光的频宽 $\Delta\nu$。
\item 若将此镉灯作为迈克耳孙干涉仪的光源，用镜面移动来观测干涉场输出光电讯号曲线，设镜面移动速度为 $0.5\,\mathrm{mm/s}$，试估算约需多长时间 $\Delta t$，可以获得显示有两个波包形状的讯号曲线。
\end{enumerate}
}
\problem{4-23}{
设有两条光谱线其波长约为 $600\,\mathrm{nm}$ 而波长差约在 $10^{-4}\,\mathrm{nm}$ 量级，现要求用法布里--珀罗干涉仪将它俩分辨开来，试问法--珀仪的镜面间距 $h$ 至少要多长？设单镜面反射率 $R$ 为 $0.95$。
}
\problem{4-24}{
设法--珀仪两镜面之间距离 $h$ 为 $1\,\mathrm{cm}$，用波长为 $500\,\mathrm{nm}$ 的绿光作实验，干涉图样的中心恰好是一亮斑，试求出第 $10$ 个亮环的角直径 $\Delta\theta$；若用一长焦距为 $300\,\mathrm{mm}$ 的镜头拍摄，所得该环的直径 $d$ 为多少？
}
\problem{4-25}{
设一法--珀腔其长度为 $5\,\mathrm{cm}$，若用准单色扩展光源作实验，其中心光波长为 $600\,\mathrm{nm}$。
\begin{enumerate}
\item 求中心干涉级数 $k_0$。
\item 在倾角为 $1^\circ$ 附近，干涉环的半角宽度 $\Delta\theta_h$ 为多少？设反射率 $R$ 为 $0.98$。
\item 若用这个法--珀腔来分辨谱线，其色分辨本领 $R_{\mathrm e}$ 为多少？可分辨的最小波长间隔 $\delta\lambda_{\mathrm m}$ 为多少？设 $\lambda\approx600\,\mathrm{nm}$。
\item 若用一束白光正入射于这法--珀腔，以使法--珀腔对白光进行选频。试问输出纵模的频率间隔 $\Delta\nu$ 为多少？其单模线宽 $\Delta\nu_h$ 为多少？这相当于谱线宽度 $\Delta\lambda_h$ 为多少？
\item 由于测量过程中的热胀冷缩，引起该法--珀腔长的改变量 $\delta h/h$ 约为 $10^{-5}$，则输出谱线的漂移量 $\delta\lambda$ 为多少？
\end{enumerate}
}
\problem{4-26}{
利用多光束干涉可以制成一种干涉滤光片，如图所示，在玻璃平晶上镀一层银，在银面上蒸镀一层透明膜，在膜上再镀一层银，于是，两个高反射率银面之间形成一个膜层而产生多光束干涉。那透明膜层材料可选用冰晶石 $(\mathrm{3NaF}\cdot\mathrm{AlF_3})$，其折射率为 $1.55$。设银面反射率 $R$ 为 $0.96$，膜层厚度 $h$ 为 $0.40\,\mu\mathrm m$。
\begin{enumerate}
\item 在可见光范围内，透射光最强的谱线有几条，其光波长各为多少？
\item 每条透射谱线的宽度 $\Delta\lambda_h$ 为多少？
\end{enumerate}
\tikzfig{fig_4_26}{4-26}
}
\problem{4-27}{
参见题图，对于完全消反射的单层膜，其光学厚度 $nh$ 和折射率 $n$，需要同时满足以下两个条件：
\begin{enumerate}
\item $nh=(2k+1)\dfrac{\lambda_0}{4}$，$k=0,1,2,3,\ldots$；
\item $n=\sqrt{n_1\cdot n_{\mathrm g}}$。
\end{enumerate}
试对此给出证明。
\tikzfig{fig_4_27}{4-27}
}
